\documentclass[conference]{IEEEtran}
\IEEEoverridecommandlockouts
\usepackage[utf8]{inputenc}
\usepackage[T1]{fontenc}
\usepackage{cite}
\usepackage{amsmath,amssymb,amsfonts}
\usepackage{graphicx}
\usepackage{textcomp}
\usepackage{xcolor}
\usepackage{booktabs}
\usepackage{array}
\usepackage{hyperref}

\def\BibTeX{{\rm B\kern-.05em{\sc i\kern-.025em b}\kern-.08em
    T\kern-.1667em\lower.7ex\hbox{E}\kern-.125emX}}

\begin{document}

\title{Athletix Governance: Design and Development of a Federal Digital Platform for the Tunisian Swimming Federation}

\author{
\IEEEauthorblockN{Alaa Ben Chouikha}
\IEEEauthorblockA{Class 1ALINFO1 \\ ESPRIT School of Engineering \\ Tunis, Tunisia \\ alaa.benchouikha@esprit.tn}
\and
\IEEEauthorblockN{Manar Ben Kacem}
\IEEEauthorblockA{Class 1ALINFO1 \\ ESPRIT School of Engineering \\ Tunis, Tunisia \\ manar.benkacem@esprit.tn}
\and
\IEEEauthorblockN{Elaa}
\IEEEauthorblockA{Class 1ALINFO1 \\ ESPRIT School of Engineering \\ Tunis, Tunisia \\ elaa@esprit.tn}
\and
\IEEEauthorblockN{Mazen Jebali}
\IEEEauthorblockA{Class 1ALINFO1 \\ ESPRIT School of Engineering \\ Tunis, Tunisia \\ mazen.jebali@esprit.tn}
\and
\IEEEauthorblockN{Montassar Tayari}
\IEEEauthorblockA{Class 1ALINFO1 \\ ESPRIT School of Engineering \\ Tunis, Tunisia \\ montassar.tayari@esprit.tn}
}

\maketitle

\begin{abstract}
The Tunisian Swimming Federation (FTN) governs all aquatic sports in Tunisia (competitive swimming, open water, artistic swimming, diving, water polo), yet currently relies on a static, content-only WordPress website with no advanced user interaction. This paper presents the study, design, and development of \emph{Athletix Governance}, a federal digital platform intended to replace this legacy site. Building on an analysis of the existing system and the needs of four identified actor categories (public, athletes, coaches/clubs, federation administrators), we specify ten functional requirements and eight non-functional requirements, then propose an eight-module architecture: six core modules (authentication, competitions, results and rankings, clubs, content, training programs) and two value-added modules (affiliated swimming-pool management with online slot booking and club lane assignment, and a community forum), the latter two being absent from every African aquatic federation website surveyed. The solution is implemented with a Spring Boot REST API, schema migrations versioned through Liquibase, JWT-based authentication, and a modular Angular frontend built on reactive signals. Results show that this architecture modernizes competition, license, and results management while remaining maintainable and scalable.
\end{abstract}

\begin{IEEEkeywords}
federal digital platform, sports competition management, Spring Boot, Angular, digital transformation, swimming federation, REST architecture
\end{IEEEkeywords}

\section{Introduction}
The Tunisian Swimming Federation (FTN) is the official body governing aquatic sports in Tunisia. It federates all clubs practicing competitive swimming, open-water swimming, artistic swimming, diving, and water polo, organizes national competitions, and manages the national teams. It is a member of \emph{Africa Aquatics} and \emph{World Aquatics}, and is headquartered at the Maison des Fédérations Sportives, Cité Olympique El Menzah, Tunis.

Despite this central role in the national sports ecosystem, the FTN's digital presence relies on an institutional website (\texttt{ftnatation.tn}) designed as a simple content showcase: news, results, calendars, and official documents are published manually, with no personal space for athletes or clubs and no transactional features (online competition registration, dematerialized license requests). This creates a growing gap between the expectations of modern digital sport users and the federation's actual digital offering.

The objective of this work is twofold: (1) conduct a structured study of the existing system to precisely identify the current site's limitations and the needs of the federation's various actors, and (2) design and develop a federal digital platform, \emph{Athletix Governance}, that addresses these needs through a modern, secure, and scalable modular architecture. The remainder of this paper is organized as follows: Section II presents the existing-system study, Section III details the actors and functional/non-functional requirements, Section IV describes the modular architecture of the solution, Section V presents the technical implementation, Section VI discusses the solution's added value, and Section VII concludes with perspectives for future work.

\section{Existing System Analysis}
\subsection{Analysis of the current website}
The website \texttt{ftnatation.tn} is built on WordPress and organizes its content into seven main sections: \emph{The Tunisian Federation} (regulations, organizational chart, affiliated clubs, technical staff), \emph{Swimming}, \emph{Open Water}, and \emph{Water Polo} (results, rankings, schedules, national teams, and national calendar for each discipline), \emph{Swimming Academy} (program and news), \emph{Licenses} (per-club listing), and \emph{Press Corner} (articles and press releases). This structure repeats the same sub-sections for every discipline without offering a consolidated view.

\subsection{Identified limitations}
The analysis reveals seven major categories of limitations.

\begin{enumerate}
\item \textbf{Aging, non-responsive interface} — outdated WordPress theme, menus and tables illegible on mobile, no UX/UI optimization since the site's creation.
\item \textbf{Inconsistent bilingual content} — no language switcher, pages either entirely in Arabic or entirely in French with no systematic translation.
\item \textbf{Lack of interactive features} — no online competition registration nor dematerialized license requests; all procedures require a phone call or a physical visit.
\item \textbf{No real-time competition management} — results published manually as PDF or HTML files, no automatic ranking updates, information often published several days late.
\item \textbf{Unintuitive navigation} — nested menus, redundant sub-sections, no internal search engine.
\item \textbf{No secure member space} — no authentication system: an athlete cannot view their performance history, a club cannot manage its members online.
\item \textbf{Rudimentary media management and no public API} — no photo/video gallery, uncategorized news, no programmatic access to federation data, technically isolating the federation from its digital ecosystem.
\end{enumerate}

The current site therefore remains a content-oriented institutional site with no advanced user interaction. This observation is the starting point for the target-solution specification presented in the following sections.

\section{Actors and Requirements}
\subsection{System actors}
Four actor categories were identified, summarized in Table~\ref{tab:actors}.

\begin{table}[h]
\caption{System actors}
\label{tab:actors}
\centering
\begin{tabular}{p{1.6cm}p{4.8cm}p{1.3cm}}
\toprule
\textbf{Actor} & \textbf{Description} & \textbf{Access} \\
\midrule
Public / Media & Unauthenticated visitors consulting results, news, calendar & Public \\
Athlete & Licensed member managing their profile, results, and registrations & Authenticated \\
Coach / Club & Responsible for managing club athletes and competition entries & Authenticated \\
Federation Admin & Platform manager (content, licenses, competitions, users) & Admin \\
\bottomrule
\end{tabular}
\end{table}

\subsection{Functional requirements}
Table~\ref{tab:fr} summarizes the ten functional requirements (FR) identified from the actor and use-case analysis.

\begin{table}[h]
\caption{Functional requirements}
\label{tab:fr}
\centering
\begin{tabular}{p{0.5cm}p{6cm}p{0.9cm}}
\toprule
\textbf{ID} & \textbf{Functional requirement} & \textbf{Prio.} \\
\midrule
FR-01 & User management: secure authentication, roles, access rights & High \\
FR-02 & Athlete management: profiles, performance history, ranking & High \\
FR-03 & Competition management: creation, calendar, results publication & High \\
FR-04 & Online dematerialized registrations and licenses & High \\
FR-05 & Event calendar with multi-criteria filters & High \\
FR-06 & Advanced search engine (athletes, competitions, clubs) & Medium \\
FR-07 & CMS for federal news and announcements & Medium \\
FR-08 & Performance dashboards and statistics & Medium \\
FR-09 & Affiliated pool management, interactive map, booking & High \\
FR-10 & Community forum with moderation and Q\&A space & Medium \\
\bottomrule
\end{tabular}
\end{table}

\subsection{Non-functional requirements}
Eight non-functional requirements frame the design: \emph{performance} (under 2 seconds load time for main pages), \emph{security} (JWT authentication, HTTPS, BCrypt hashing, GDPR compliance), \emph{responsive design} (mobile, tablet, desktop adaptation), \emph{accessibility} (WCAG 2.1 level AA compliance), \emph{multilingual support} (Arabic RTL, French, English), \emph{scalability} (load spikes during major competitions), \emph{maintainability} (modular, documented, tested code), and \emph{reliability} (minimum 99.5\% availability).

\section{Solution Architecture}
The \emph{Athletix Governance} solution is divided into eight independent functional modules, each covering a complete business domain from the backend layer to the user interface. The first six modules are \emph{core} modules that modernize functionality already present on the existing site; the last two are \emph{value-added} modules, introducing services absent from the current site and, to our knowledge, from any other African aquatic federation website.

\subsection{Module 1 — Authentication and User Management}
Secure JWT-based authentication and role management, personal profile editing, license request submission and validation, full user CRUD with access-rights assignment.

\subsection{Module 2 — Competition Management}
Competition creation (discipline, dates, venue, categories), calendar browsing with filters (discipline, region, level), individual registration with automatic eligibility validation, bulk registration by coaches, modification/cancellation with automatic notification.

\subsection{Module 3 — Results and Rankings}
Official results entry and publication after each competition, filtered consultation by type, category, year, gender, and discipline, personal performance history, national ranking by discipline and age category, club performance dashboard.

\subsection{Module 4 — Club Management}
Public athlete profiles (record, club, discipline, nationality), club member CRUD with license tracking and competition assignment, affiliated club directory, key indicators (active licenses, clubs and athletes per region), assignment of a training pool and lane to each club, visible on its public page.

\subsection{Module 5 — Content and Backoffice}
Categorized news articles with pagination and search, CMS for publishing announcements and official documents, unified global search (athletes, competitions, clubs, articles), admin-managed media gallery, public dashboard of aggregate statistics (licenses, clubs, athletes) feeding the home page.

\subsection{Module 6 — Training Programs}
This module digitizes the \emph{Swimming Academy} section of the existing site: catalogue of training programs by age bracket (description, minimum/maximum age, visual, active/inactive status), full CRUD on the admin side, public listing of active programs on the home page to guide families toward the program suited to a child's age.

\subsection{Module 7 — Pool Management (Value-Added)}
This module centralizes affiliated aquatic infrastructure and introduces an online slot-booking system: geolocated interactive map, detailed per-pool sheet (capacity, 25m/50m pool types, opening hours, equipment), training slot booking, availability management and unavailability reporting by the administrator, multi-criteria filtering. It is also the source of the lanes assigned to clubs for their training sessions (see Module 4).

\subsection{Module 8 — Community Forum (Value-Added)}
A discussion space structured by category (swimming, water polo, open water, training), thread creation and replies, social interactions (like, report), pinning of official announcements, admin moderation, and a dedicated Q\&A space between coaches and athletes.

Table~\ref{tab:modules} summarizes the use-case distribution per module.

\begin{table}[h]
\caption{Module summary}
\label{tab:modules}
\centering
\begin{tabular}{p{0.9cm}p{3.3cm}p{1cm}p{1.6cm}}
\toprule
\textbf{Module} & \textbf{Name} & \textbf{UC} & \textbf{Type} \\
\midrule
M1 & Authentication \& Users & 5 & Core \\
M2 & Competition Management & 5 & Core \\
M3 & Results \& Rankings & 5 & Core \\
M4 & Athletes \& Clubs & 4 & Core \\
M5 & Content \& Backoffice & 4 & Core \\
M6 & Training Programs & 4 & Core \\
M7 & Pool Management & 6 & Value-added \\
M8 & Community Forum & 7 & Value-added \\
\bottomrule
\end{tabular}
\end{table}

\section{Technical Implementation}
\subsection{Backend}
The backend exposes a REST API built with \emph{Spring Boot}, organized into \texttt{Controller}/\texttt{Service}/\texttt{Repository} layers per business domain (competitions, events, results, rankings, club events, clubs, athletes, news). Persistence relies on a relational database whose schema is version-controlled with \emph{Liquibase}: every schema change is described by an XML \emph{changeset}, and table-creation changesets are made idempotent through \texttt{preConditions} (\texttt{onFail="MARK\_RAN"}) checking for the prior absence of the targeted table or column. This allows migrations originating from parallel development branches to be merged safely without risking a conflict on an already-initialized database. Authentication relies on JWT tokens with role-based access control (RBAC), in line with the security non-functional requirement specified in Section III-C.

\subsection{Frontend}
The frontend is built with \emph{Angular}, using \emph{standalone} components organized into lazily-loaded feature modules for each functional module (competitions, results, athletes/clubs, content, pools, forum, administration). Reactive component state is managed through the \emph{signals} API (\texttt{signal}, \texttt{computed}) rather than classic observable streams, simplifying local state management (filters, pagination, search) and improving code readability.

\subsection{Deployment}
The entire stack (backend API, database, frontend) is containerized and orchestrated through \emph{Docker Compose}, enabling reproducible deployment on any environment with a Docker engine.

\section{Discussion}
Modules 7 (pool management) and 8 (community forum) constitute the main added value of this platform compared to the existing site. To our knowledge, no African aquatic federation currently offers an online pool slot-booking system: this module reduces informal phone-based procedures, optimizes the use of national sports infrastructure, and constitutes a genuine differentiator for the FTN within its regional ecosystem. Similarly, the community forum centralizes exchanges currently scattered across unofficial social networks (WhatsApp, Facebook), while showcasing coach expertise through a dedicated Q\&A space and giving the federation an official, moderated communication channel.

On the technical side, the choice of a modular architecture with idempotent migrations facilitates the continuous integration of modules developed in parallel by separate teams — a concrete constraint encountered during development, where several feature branches (competitions, news) had to be merged into the integration branch without regressing the database schema.

\section{Conclusion and Future Work}
This paper presented the existing-system study, requirements specification, and architecture of the \emph{Athletix Governance} federal digital platform, designed to replace the Tunisian Swimming Federation's current institutional website. The analysis revealed a limited site with no interaction or personal space for its actors; the proposed solution addresses these limitations through six core modules that modernize existing functionality and two value-added modules (pools, forum) that position the platform as a pioneer within the African sports federation ecosystem. The technical implementation carried out — Spring Boot REST API, idempotent Liquibase migrations, JWT authentication, reactive Angular frontend — demonstrates the viability of this architecture.

Future work includes full interface internationalization (Arabic RTL, French, English), deployment of a companion mobile application, and the introduction of real-time notifications (results, license validation, pool slot availability) via WebSocket.

\section*{Acknowledgment}
The authors thank ESPRIT School of Engineering and the Tunisian Swimming Federation for the business context that made this academic project possible.

\begin{thebibliography}{00}
\bibitem{springboot} VMware/Spring Team, ``Spring Boot Reference Documentation,'' [Online]. Available: https://docs.spring.io/spring-boot/
\bibitem{angular} Google, ``Angular Documentation,'' [Online]. Available: https://angular.dev
\bibitem{jwt} M. Jones, J. Bradley, N. Sakimura, ``JSON Web Token (JWT),'' RFC~7519, IETF, May 2015.
\bibitem{owasp} OWASP Foundation, ``OWASP Top Ten,'' [Online]. Available: https://owasp.org/www-project-top-ten/
\bibitem{wcag} W3C, ``Web Content Accessibility Guidelines (WCAG) 2.1,'' [Online]. Available: https://www.w3.org/TR/WCAG21/
\bibitem{worldaquatics} World Aquatics, Official Website, [Online]. Available: https://www.worldaquatics.com
\bibitem{africaaquatics} Africa Aquatics, Official Website, [Online]. Available: https://africaaquatics.org
\bibitem{liquibase} Liquibase Inc., ``Liquibase Documentation,'' [Online]. Available: https://docs.liquibase.com
\bibitem{postgresql} PostgreSQL Global Development Group, ``PostgreSQL Documentation,'' [Online]. Available: https://www.postgresql.org/docs/
\end{thebibliography}

\end{document}
